% META: {"id": "TCOMPL-BAYES-EX-017", "chapitre": "TCOMPL-INFERENCE-BAYESIENNE", "type_objet": "exercice", "capacites": ["TCOMPL-BAYES-C2"], "capacites_codes": ["C2"], "methodes": ["M2"], "parcours": 2, "competences": ["communiquer","raisonner"], "duree_min": 20, "mode_creation": "ex_nihilo", "sources_inspiration": [], "fichier_tex": "chapitres/TCOMPL-INFERENCE-BAYESIENNE/exercices/TCOMPL-BAYES-EX-017.tex", "corrige_tex": "chapitres/TCOMPL-INFERENCE-BAYESIENNE/corriges/TCOMPL-BAYES-CO-017.tex", "status": "approved"}
% BEGIN-VERIFY
% from sympy import *
% x = symbols('x')
% f = x**3 - 6*x**2 + 9*x + 1
% fp = diff(f, x)
% assert expand(fp) == 3*x**2 - 12*x + 9
% assert solve(fp, x) == [1, 3]
% assert f.subs(x, 1) == 5
% assert f.subs(x, 3) == 1
% # tangente en x=2: f(2)=3, f'(2)=-3, y=-3(x-2)+3=-3x+9
% assert f.subs(x, 2) == 3
% assert fp.subs(x, 2) == -3
% END-VERIFY
\begin{exercice}{TCOMPL-BAYES-EX-017}{2}{20}
Soit $f(x) = x^3 - 6x^2 + 9x + 1$.
\begin{enumerate}
  \item Dresser le tableau de variations de $f$.
  \item Determiner l'equation de la tangente a la courbe au point d'abscisse $2$.
  \item La courbe traverse-t-elle cette tangente en $x = 2$ ?
\end{enumerate}

\end{exercice}
